\section*{Chapitre \ref{ch:b1:comparative-planetology} --- Planétologie comparée}

\begin{solution}{exo:b1:comparative-planetology:1}
$V = \tfrac43\pi(\qty{6.9911e7}{m})^3 = \qty{1.4313e24}{m^3}$ et
$\bar\rho = \num{1898.13e24}/\num{1.4313e24} = \qty{1326}{kg.m^{-3}}$, soit
seulement 1.33 fois la masse volumique de l'eau, bien que Jupiter soit 318 fois
plus massive que la Terre~: elle est faite surtout des éléments les plus légers.
\end{solution}

\begin{solution}{exo:b1:comparative-planetology:2}
Mercure (rayon \qty{2440}{km}, \qty{5429}{kg.m^{-3}}) et Vénus
(\qty{6052}{km}, \qty{5243}{kg.m^{-3}}) sont petites et plus denses que la
roche~: \omterm{def:b1:comparative-planetology:planet-types}{planètes telluriques}, de roche et de fer. Saturne (\qty{60268}{km},
\qty{687}{kg.m^{-3}}) est énorme et plus légère que l'eau~: une \omterm{def:b1:comparative-planetology:planet-types}{géante gazeuse},
faite surtout d'hydrogène et d'hélium. Uranus (\qty{25559}{km},
\qty{1270}{kg.m^{-3}}) est une géante mais bien plus petite que Saturne et
presque deux fois plus dense~: une \omterm{def:b1:comparative-planetology:planet-types}{géante de glaces}, dont l'eau, le méthane et
l'ammoniac sont plus lourds que l'hydrogène.
\end{solution}

\begin{solution}{exo:b1:comparative-planetology:3}
$T_\mathrm{eq} = (2601.3\times0.23/(4\times\num{5.6704e-8}))^{1/4} =
(\num{2.638e9})^{1/4} = \qty{226.6}{K}$. Vénus reçoit 1.9 fois plus de lumière
solaire que la Terre mais ses nuages en réfléchissent \qty{77}{\percent}~:
elle n'absorbe que $2601.3\times0.23 = \qty{598}{W.m^{-2}}$ contre
$1361\times0.706 = \qty{961}{W.m^{-2}}$ pour la Terre.
\end{solution}

\begin{solution}{exo:b1:comparative-planetology:4}
De la plus homogène à la plus concentrée~: la Lune, Mars, la Terre, Jupiter,
Saturne. Une masse ponctuelle au centre a $I = 0$, donc $C = 0$.
\end{solution}

\begin{solution}{exo:b1:comparative-planetology:5}
$T = 15.945421\times86400 = \qty{1377684}{s}$~;
$GM = 4\pi^2(\num{1.22187e9})^3/\num{1377684}^2 = \qty{3.794e16}{m^3.s^{-2}}$~;
$M = \num{3.794e16}/\num{6.6743e-11} = \qty{5.685e26}{kg}$, à moins de
\qty{0.03}{\percent} de \qty{568.32e24}{kg}.
\end{solution}

\begin{solution}{exo:b1:comparative-planetology:6}
$S = 1361\times(149.6/228.0)^2 = \qty{585.9}{W.m^{-2}}$ à Mars et
$1361\times(149.6/778.5)^2 = \qty{50.26}{W.m^{-2}}$ à Jupiter, en accord
avec les tables (la petite différence à Mars vient de l'arrondi des distances
moyennes).
\end{solution}

\begin{solution}{exo:b1:comparative-planetology:7}
$S \propto d^{-2}$ et $T_\mathrm{eq} \propto S^{1/4}$, donc $T_\mathrm{eq}
\propto d^{-1/2}$. Pour Jupiter, $T_\mathrm{eq} = (50.26\times0.657/(4\times
\num{5.6704e-8}))^{1/4} = \qty{109.8}{K}$. Une planète qui rayonne plus qu'elle
n'absorbe doit avoir une source interne de chaleur~: Jupiter libérerait de la
chaleur stockée depuis sa formation, ou produite par sa lente contraction.
\end{solution}

\begin{solution}{exo:b1:comparative-planetology:8}
Vénus~: $737 - 226.6 = \qty{510}{K}$, soit \qty{225}{\percent} de $T_\mathrm{eq}$.
La Terre~: $T_\mathrm{eq} = (1361\times0.706/(4\sigma))^{1/4} = \qty{255.1}{K}$,
$288 - 255.1 = \qty{33}{K}$, soit \qty{13}{\percent}. Mars~: $214 - 209.8 =
\qty{4}{K}$, soit \qty{2}{\percent}.
\end{solution}

\begin{solution}{exo:b1:comparative-planetology:9}
Avec $\mathrm dm = 4\pi\rho r^2\,\mathrm dr$~: $I = \int_0^R\tfrac23r^2\,4\pi
\rho r^2\,\mathrm dr = \tfrac{8\pi}{15}\rho R^5$, et $MR^2 = \tfrac43\pi\rho
R^5$, donc $C = \tfrac{8}{15}\cdot\tfrac34 = \tfrac25$. Une coquille creuse a
toute sa masse à $r = R$~: $I = \tfrac23MR^2$, $C = \tfrac23$.
\end{solution}

\begin{solution}{exo:b1:comparative-planetology:10}
$\tfrac52C\bar\rho = 2.5\times0.35\times5429 = 4750.4$~;
$x^2 = (4750.4 - 3300)/(5429 - 3300) = 1450.4/2129 = 0.6812$, $x = 0.8254$,
rayon du noyau $0.8254\times2439.7 = \qty{2014}{km}$. $x^3 = 0.5623$ et
$\rho_c = 3300 + 2129/0.5623 = \qty{7086}{kg.m^{-3}}$. Le noyau contient
$0.5623\times7086/5429 = \qty{73}{\percent}$ de la masse. Le noyau de Mercure
atteint \qty{83}{\percent} de son rayon, contre $3485/6371 =
\qty{55}{\percent}$ pour la Terre~: Mercure est presque tout noyau, avec un
mince manteau rocheux.
\end{solution}

\begin{solution}{exo:b1:comparative-planetology:11}
(a) La fonction $t\mapsto t^4$ est convexe, donc
$\bigl(\tfrac{T_1 + T_2}{2}\bigr)^4 \le \tfrac{T_1^4 + T_2^4}{2} =
T_\mathrm{eq}^4$, avec égalité seulement si $T_1 = T_2$. (b) $T_1^4 =
2\times270^4 - 100^4 = \num{1.0529e10}$, $T_1 = \qty{320}{K}$, et la température
moyenne est $(320 + 100)/2 = \qty{210}{K}$, soit \qty{60}{K} sous
$T_\mathrm{eq}$.
\end{solution}

\begin{solution}{exo:b1:comparative-planetology:12}
$m = 4\pi R^2p/g = 4\pi\times(\num{6.0518e6})^2\times\num{9.2e6}/8.9 =
\num{4.602e14}\times\num{1.034e6} = \qty{4.76e20}{kg}$~: 93 fois la masse de
l'atmosphère terrestre, et \num{9.8e-5} de la masse de Vénus. La Terre contient
une quantité comparable de carbone, non dans son air mais dans les roches
carbonatées (calcaires) et dans l'océan, où l'eau liquide a permis de le stocker.
\end{solution}

\begin{solution}{pb:b1:comparative-planetology:1}
\textbf{1.} $T = 0.31891\times86400 = \qty{27553.8}{s}$.

\textbf{2.} $GM = 4\pi^2\times(\num{9.378e6})^3/\num{27553.8}^2 =
\qty{4.289e13}{m^3.s^{-2}}$.

\textbf{3.} $M = \num{4.289e13}/\num{6.6743e-11} = \qty{6.426e23}{kg}$,
à \qty{0.14}{\percent} de la valeur tabulée \qty{6.4169e23}{kg}.

\textbf{4.} $\num{1.06e16}/\num{6.426e23} = \num{1.6e-8}$~: la masse de Phobos
est négligeable.

\textbf{5.} $v = 2\pi a/T = 2\pi\times\num{9.378e6}/27553.8 =
\qty{2138}{m.s^{-1}}$, soit environ \qty{2.1}{km.s^{-1}}.

\textbf{6.} $T = \qty{7.65}{h}$, 3.2 fois plus courte que la rotation de Mars~:
Phobos tourne plus vite que le sol ne tourne sous lui, il dépasse donc la
rotation. Vu de Mars, il se lève à l'ouest et se couche à l'est.

\textbf{7.} $V = \tfrac43\pi(\num{3.3895e6})^3 = \qty{1.6312e20}{m^3}$.

\textbf{8.} $\bar\rho = \num{6.426e23}/\num{1.6312e20} = 3939$, soit
\qty{3940}{kg.m^{-3}}.

\textbf{9.} $3940/5514 = 0.715$. La Terre, plus grande et plus massive,
comprime davantage son propre intérieur~; et son noyau de fer est une plus
grande part de la planète.

\textbf{10.} $g = GM/R^2 = \num{4.289e13}/(\num{3.3895e6})^2 =
\qty{3.73}{m.s^{-2}}$, soit \qty{38}{\percent} de celle de la Terre.

\textbf{11.} $v_\mathrm{esc} = \sqrt{2GM/R} = \sqrt{2\times\num{4.289e13}/
\num{3.3895e6}} = \qty{5.03}{km.s^{-1}}$.

\textbf{12.} $T_\mathrm{eq} = (586.2\times0.75/(4\times\num{5.6704e-8}))^{1/4}
= \qty{209.8}{K}$.

\textbf{13.} $214 - 209.8 = \qty{4}{K}$~: un faible réchauffement par effet de serre.

\textbf{14.} $m = p/g = 636/3.733 = \qty{170}{kg.m^{-2}}$~; pour la Terre
$101400/9.82 = \qty{10326}{kg.m^{-2}}$, soit 61 fois plus.

\textbf{15.} $4\pi R^2\times170.4 = \num{1.444e14}\times170.4 =
\qty{2.46e16}{kg}$.

\textbf{16.} Au-dessus de chaque mètre carré, il y a soixante fois moins de gaz
que sur la Terre, si bien que l'absorption de l'infrarouge, bien que due au
dioxyde de carbone, est faible~; et l'air de Mars contient très peu de vapeur
d'eau, l'autre grand absorbant.

\textbf{17.} Une Mars homogène aurait $C = 0.4$~; la valeur mesurée, 0.366,
montre que sa masse volumique croît vers le centre~: Mars a un noyau dense.

\textbf{18.} $\bar\rho = \rho_m + (\rho_c - \rho_m)x^3$.

\textbf{19.} $C = \tfrac25[\rho_m + (\rho_c - \rho_m)x^5]/\bar\rho$.

\textbf{20.} D'après 18 et 19, $(\rho_c - \rho_m)x^3 = \bar\rho - \rho_m$ et
$(\rho_c - \rho_m)x^5 = \tfrac52C\bar\rho - \rho_m$~; en divisant, on obtient $x^2$.

\textbf{21.} $\tfrac52C\bar\rho = 2.5\times0.366\times3940 = 3605.1$~;
$x^2 = (3605.1 - 3500)/(3940 - 3500) = 105.1/440 = 0.2389$, $x = 0.4887$.
$x^3 = 0.1167$ et $\rho_c = 3500 + 440/0.1167 = \qty{7269}{kg.m^{-3}}$.

\textbf{22.} $x^3\rho_c/\bar\rho = 0.1167\times7269/3940 = 0.215$, soit
\qty{21.5}{\percent} de la masse.

\textbf{23.} Avec $\rho_m = 3450$~: $x^2 = 155.1/490 = 0.3165$, $x = 0.5626$,
rayon du noyau \qty{1907}{km}, $\rho_c = \qty{6202}{kg.m^{-3}}$. Avec $\rho_m =
3550$~: $x^2 = 55.1/390 = 0.1413$, $x = 0.3759$, \qty{1274}{km},
$\rho_c = \qty{10894}{kg.m^{-3}}$. Une variation de \qty{1.4}{\percent} de la
masse volumique supposée du manteau déplace le rayon du noyau de plusieurs
centaines de kilomètres~: la masse volumique moyenne et $C$ seuls ne peuvent
fixer précisément le noyau de Mars, ce qui explique l'envoi de sismomètres.

\textbf{24.} $0.4887\times3389.5 = \qty{1657}{km}$~: le rayon du noyau de Mars
d'après le modèle est de \textbf{environ \qty[text-series-to-math=true]{1660}{km}}, dans l'intervalle
sismique récent de \qty{1650}{km} à \qty{1675}{km} et bien en dessous des
\qty{1830}{km} du tableau planétaire --- un accord que la question 23 montre
être en partie dû à la chance dans le choix de $\rho_m$.
\end{solution}
